% =============================================================================
% Papers/prompt-to-render-educational-video-generation.tex
%
% Prompt to Render: A Systematic Literature Review on Large Language Models
% for Educational Video/Animation Generation
%
% Standalone research survey (IEEE journal format). Not part of the CS401
% coursework tree -- lives in Papers/, independent of Slides/Quarto/Notes.
%
% Source material: quality_reports/lit_review_prompt_to_render_educational_
% video_generation.md -- 24 primary studies (Round 1 + Round 2), plus 2
% methodology references and 3 Related-Work secondary studies, all
% independently verified via forked claim-verifier passes (CoVe) before this
% manuscript was drafted or expanded from them. Round 1 (2026-08-06) produced
% a narrative survey of 15 sources; Round 2 (2026-08-13) added formal SLR
% machinery (research questions, search strategy, inclusion/exclusion
% criteria, PRISMA-style selection accounting), an animation-specific
% literature thread, and a wider search for human-subject learning-outcome
% studies, per the plan at quality_reports/plans/eager-crunching-magpie.md.
%
% Compile with:
%   cd Papers
%   xelatex -interaction=nonstopmode prompt-to-render-educational-video-generation.tex
%   bibtex prompt-to-render-educational-video-generation
%   xelatex -interaction=nonstopmode prompt-to-render-educational-video-generation.tex
%   xelatex -interaction=nonstopmode prompt-to-render-educational-video-generation.tex
% =============================================================================

\documentclass[journal]{IEEEtran}
\usepackage{amsmath,amssymb}
\usepackage{graphicx}
\usepackage{hyperref}
\hypersetup{colorlinks=true, linkcolor=blue, citecolor=blue, urlcolor=blue}

\begin{document}

\title{Prompt to Render: A Systematic Literature Review on Large Language Models for Educational Video/Animation Generation}

\author{
\IEEEauthorblockN{Dr. Auqib Hamid Lone}
\IEEEauthorblockA{Department of Computer Science \& Engineering\\
Government College of Engineering \& Technology (GCET), Ganderbal\\
% Email intentionally left as a placeholder -- fill in before submission.
Email: \texttt{[your-email]@gcetganderbal.ac.in}}
}

\maketitle

\begin{abstract}
Large language models are increasingly used not only to draft educational text but to orchestrate the generation of entire instructional videos, narrated animations, and slide decks. This systematic literature review (SLR) examines that shift by following an explicit, reported methodology -- research questions, a documented search strategy, inclusion/exclusion criteria, and a PRISMA-style accounting of the study-selection process -- rather than an ad hoc narrative survey. From an initial pool of candidate records identified through web-based search of arXiv, ACL Anthology, CVF Open Access, and general indexed sources, we screened 32 new candidate records in a second search round (added to a first round conducted one week earlier) and arrived at a final corpus of 24 primary studies -- systems, benchmarks, and human-subject evaluations -- supplemented by three prior secondary studies discussed for positioning and two methodology references. We organize the corpus around three questions: the \emph{architectures} used to go from a prompt to a rendered artifact, the \emph{benchmarks} used to compare systems, and the \emph{gaps} between what current evaluation measures and what pedagogical effectiveness actually requires. We find that the field has moved decisively away from end-to-end, pixel-space text-to-video generation toward agentic pipelines that decompose the task into checkable intermediates -- executable scripts, code, verifiable traces, stroke sequences, or editable slides -- specifically because raw video diffusion offers no mechanism to guarantee factual or logical correctness. This decomposition principle now extends to a previously under-covered animation-specific thread (Manim-code generation, algorithm-visualization traces, whiteboard/stroke synthesis) that this review adds to the literature map. Benchmarking has bifurcated accordingly: mature, general-purpose video-quality suites (VBench, EvalCrafter) coexist with a growing set of education-specific benchmarks (SlidesBench, MMMC/TeachQuiz, the Paper2Video benchmark, ManimBench, EduRequire-500, TheoremExplainBench) that attempt to measure knowledge transfer or layout correctness rather than visual fidelity alone. The central finding of this review is that, of the 24 primary studies, only three report a controlled human-subject study of learning outcomes and only one additional study reports a real classroom deployment with student-perceived quality judgments; the remainder evaluate the generated artifact, not the learner. This is a meaningful increase over the two such studies (one controlled learning-outcome study, one classroom deployment) identified in a narrower first-round search, but the underlying imbalance -- roughly one in eight studies measuring a learning outcome -- persists even as the field's technical sophistication grows quickly. We argue this evaluation gap is the field's most consequential open problem and outline concrete next steps -- cross-system human-subject studies, validation of LLM-as-judge knowledge-transfer proxies against human performance, dedicated factuality and accessibility metrics, and a shared benchmark analogous to VBench but for pedagogical effectiveness -- that would close it.
\end{abstract}

\begin{IEEEkeywords}
systematic literature review, educational video generation, large language models, multi-agent systems, animation generation, text-to-video, learning outcomes evaluation, video generation benchmarks
\end{IEEEkeywords}

\section{Introduction}
\label{sec:intro}

Producing an educational video or narrated animation has traditionally meant hours of slide design, scripting, recording, and editing for a handful of minutes of finished content. Over the past two years, large language models (LLMs) have been proposed as a way to compress that pipeline: given a topic, a document, or a scientific paper, an LLM-driven system now routinely produces a full instructional video, a narrated STEM animation, or a synthetic on-screen presenter. The pace of publication in this space has accelerated sharply -- the majority of the studies surveyed in this paper appeared within a twelve-month window -- and the underlying architectural philosophy has shifted just as quickly.

Early text-to-video systems treated video generation as a single, end-to-end pixel-synthesis problem: a diffusion or transformer-based model consumes a prompt and emits a video directly. That approach is well suited to open-ended, creative video generation, but it is poorly suited to education, where a video's value depends on whether its claims are correct and whether its explanation is logically sound -- properties that pixel-space generation has no mechanism to guarantee or even check. The literature surveyed here shows the field responding to this mismatch by abandoning end-to-end generation almost entirely for educational use cases, in favor of \emph{agentic} pipelines that insert a structured, human- or machine-checkable intermediate representation (a script, a graph, executable code, a verifiable trace, a stroke sequence, or an editable slide deck) between the LLM and the final rendered output.

This paper's predecessor was a narrative survey of fifteen sources. That format let us describe what the field does, but it could not answer, in a reportable and auditable way, \emph{how thoroughly} the literature had been searched or \emph{why} a given study was included or excluded. This revision addresses that gap directly: it reports an explicit search strategy, inclusion/exclusion criteria, and a PRISMA-style~\cite{page2021prisma} accounting of the study-selection process, following the general three-phase (planning, conducting, reporting) structure recommended for software-engineering SLRs by Kitchenham and Charters~\cite{kitchenham2007guidelines}. It also closes a scope gap in the original survey: the requested title names \emph{video/animation} generation, but the first search round centered almost entirely on video pipelines, with only one animation-adjacent system (Code2Video's Manim backend) represented. This revision adds a dedicated animation-specific thread -- Manim-code generation, algorithm-visualization traces, and whiteboard/stroke-sequence synthesis -- found through a second, explicitly documented search round.

This survey is organized around three research questions that we argue must be answered together, not separately, to understand the current state of the field:

\begin{enumerate}
\item \textbf{RQ1 -- Architectures} (Section~\ref{sec:architectures}): What computational pipelines are actually used to go from a prompt to a rendered educational video, animation, or presentation, and what problem does each design choice solve?
\item \textbf{RQ2 -- Benchmarks} (Section~\ref{sec:benchmarks}): What datasets and evaluation suites exist to compare these systems, and along what dimensions do they actually measure quality?
\item \textbf{RQ3 -- Evaluation gaps} (Section~\ref{sec:gaps}): Where does the field's evaluation practice diverge from what pedagogical effectiveness actually requires, and how large is that divergence in the current literature?
\end{enumerate}

Our review of 24 primary studies -- systems, benchmarks, and human-subject evaluations, published primarily in 2025--2026 -- finds a clear and consistent answer to the third question: evaluation practice in this field overwhelmingly measures properties of the generated \emph{artifact} -- visual quality, design coherence, textual similarity to a reference -- rather than properties of the \emph{learner} who watches it. Section~\ref{sec:gaps} quantifies this claim across the expanded corpus and Section~\ref{sec:discussion} argues it remains the single highest-leverage open problem for future work in this area.

\subsection{Outline}

The remainder of this paper is organized as follows. Section~\ref{sec:related} positions this SLR against three prior secondary studies that touch adjacent ground. Section~\ref{sec:methodology} describes the review methodology: research questions, search strategy, selection criteria, and the study-selection process. Section~\ref{sec:architectures} addresses RQ1 by examining generation architectures, including the new animation-specific thread. Section~\ref{sec:benchmarks} addresses RQ2 by cataloguing general-purpose and education-specific benchmarks. Section~\ref{sec:gaps} addresses RQ3 by quantifying what current evaluation practice measures versus what pedagogical effectiveness requires. Section~\ref{sec:threats} discusses threats to the validity of this review. Section~\ref{sec:discussion} discusses implications and concrete next steps. Section~\ref{sec:conclusion} concludes.

\section{Related Work}
\label{sec:related}

Three prior secondary studies touch ground adjacent to this SLR, and each differs from it in scope. Table~\ref{tab:related-work} summarizes the comparison.

Liu, Xiang, Li et al.~\cite{liu2024aigve} formally establish ``AI-Generated Video Evaluation'' (AIGVE) as a distinct research problem, taxonomizing existing approaches into metric-based, human-involved, and emerging model-centered (LLM/VLM-as-judge) evaluation. Their survey is general-purpose: it covers video generation evaluation broadly, with no education-specific framing, no coverage of generation architectures, and no discussion of learning outcomes. This SLR's evaluation-gaps analysis (Section~\ref{sec:gaps}) can be read as a narrowing of the AIGVE critique to the specific axis that survey does not focus on: pedagogical effectiveness.

Zhang and Shukor~\cite{zhang2025roleaivideos} conduct a PRISMA-based SLR screening 3{,}271 initial records across six databases down to 21 included studies, examining the role of AI-generated instructional video in higher education and the educational theories that support its use. Their review is pedagogy-theory-focused; it does not survey generation architectures, code/script/trace intermediates, or technical benchmarks -- the two pillars (architectures and benchmarks) this SLR's title promises.

Shu, Kou, Zhang et al.~\cite{shu2025aigcvideosslr} report a smaller SLR (143 initial records across four databases -- Web of Science, ERIC, IEEE Xplore, and Google Scholar -- down to 12 included studies) on the learning impacts and applications of AI-generated instructional video, concluding that generated-video quality is ``sufficient for educational use'' but that such systems ``cannot yet replace traditionally produced videos,'' citing limitations in learning motivation and social presence. Like Zhang and Shukor, this review is learning-impact-focused and does not catalogue generation architectures or technical benchmarks.

None of the three prior reviews combines architecture-level analysis, benchmark cataloguing, and an animation-specific thread with a quantified evaluation-gap analysis, which is the combination this SLR provides.

\begin{table}[t]
\caption{Comparison of prior secondary studies. Columns indicate whether each review substantively addresses the corresponding topic.}
\label{tab:related-work}
\centering
\begin{tabular}{p{1.6cm}p{0.7cm}p{1.0cm}p{1.3cm}p{1.0cm}}
\hline
\textbf{Secondary study} & \textbf{Yr.} & \textbf{Archit.} & \textbf{Bench.} & \textbf{Learn.-gap} \\
\hline
Liu et al.~\cite{liu2024aigve} (AIGVE) & 2024 & & (general) & \\
Zhang \& Shukor~\cite{zhang2025roleaivideos} & 2025 & & & \checkmark \\
Shu et al.~\cite{shu2025aigcvideosslr} & 2025 & & & \checkmark \\
\textbf{This SLR} & 2026 & \checkmark & \checkmark & \checkmark \\
\hline
\end{tabular}
\end{table}

\section{Methodology}
\label{sec:methodology}

We conducted this SLR following the general three-phase structure (planning, conducting, reporting) proposed by Kitchenham and Charters~\cite{kitchenham2007guidelines} for software-engineering literature reviews, and report the study-selection process using a PRISMA-style~\cite{page2021prisma} accounting adapted to a single-reviewer, web-search-based review (see Section~\ref{sec:threats} for the limitations this entails).

\subsection{Research Questions}

The three research questions stated in Section~\ref{sec:intro} guided this review: \textbf{RQ1} (architectures), \textbf{RQ2} (benchmarks), and \textbf{RQ3} (evaluation gaps). RQ3 has an implicit sub-question that Section~\ref{sec:gaps} answers directly: \textbf{RQ3.1} -- of the studies that generate educational video, animation, or slide content, what fraction report any human-subject evaluation, and what fraction specifically report a controlled learning-outcome study?

\subsection{Search Strategy}

The review was conducted in two rounds. \textbf{Round 1} (2026-08-06) used web-based search targeting architectures, benchmarks, and evaluation practice for LLM-based educational video generation broadly; the specific query strings used in that round were not separately preserved in the session log, only the resulting included set -- a limitation discussed in Section~\ref{sec:threats}. \textbf{Round 2} (2026-08-13) used web-based search (arXiv, ACL Anthology, CVF Open Access, and general indexed web sources, via search-engine queries rather than direct Boolean queries against subscription databases such as IEEE Xplore or ACM DL) with the following representative query strings, fully logged:

\begin{quotation}
\ttfamily\small
``LLM generated animation education Manim''\\
``automatic whiteboard explainer video generation LLM''\\
``algorithm visualization generation LLM code animation education''\\
``agentic multi-agent educational video generation system''\\
``lecture video generation avatar LLM classroom deployment''\\
``student learning outcomes AI generated video randomized controlled study''
\end{quotation}

Round 2 additionally used targeted follow-up searches to disambiguate two unrelated systems that both name themselves ``ManimAgent,'' and to identify prior secondary studies for Section~\ref{sec:related}.

\subsection{Selection Criteria}

Table~\ref{tab:selection-criteria} reports the inclusion and exclusion criteria applied during screening.

\begin{table}[t]
\caption{Inclusion and exclusion criteria for paper selection.}
\label{tab:selection-criteria}
\centering
\begin{tabular}{p{8.2cm}}
\hline
\textbf{Inclusion Criteria} \\
\hline
\textbf{IN1:} Describes an LLM- or VLM-based system that generates an educational video, animation, whiteboard drawing, or narrated-slide artifact from a prompt, document, or paper. \\
\textbf{IN2:} Introduces a benchmark, dataset, or evaluation framework specifically for such systems. \\
\textbf{IN3:} Reports a human-subject study evaluating learning outcomes, perceived quality, or classroom deployment of such generated content. \\
\textbf{Note:} A paper must satisfy IN1 or IN2 or IN3. \\
\hline
\textbf{Exclusion Criteria} \\
\hline
\textbf{EX1:} Not written in English. \\
\textbf{EX2:} No accessible full text. \\
\textbf{EX3:} General-purpose video/animation generation with no education-specific framing (retained only as brief narrative context, e.g.\ Sora, CogVideoX, ViMax -- not counted in the included corpus). \\
\textbf{EX4:} The system produces no pre-rendered artifact (e.g.\ a real-time, live-taught interactive agent). \\
\textbf{EX5:} Secondary studies (surveys/SLRs) -- discussed in Related Work (Section~\ref{sec:related}) instead of the primary corpus. \\
\textbf{EX6:} Generative-AI-in-education work with no specific video/animation/slide generation system (e.g.\ lesson-plan generation, general tutoring-agent studies). \\
\hline
\end{tabular}
\end{table}

\subsection{Study Selection Process}

Round 2 identified 32 unique candidate records via the search strategy above. After title/abstract screening against Table~\ref{tab:selection-criteria}, 10 were included and 22 were excluded, distributed across exclusion reasons as follows: 11 records under EX4/EX6 (live agents, lesson plans, general tutoring-effectiveness studies, or gray literature such as blog posts), 8 records under EX3 (general-purpose, non-education-specific video/animation/document generation), 2 records excluded as thematically redundant with an already-included system in the same architectural thread, and 1 record excluded as pre-dating the LLM-based scope of this review. Round 1 (2026-08-06) used the same general web-search methodology and contributed 15 records to the included set, of which one (the AIGVE meta-survey~\cite{liu2024aigve}) is reclassified under EX5 into Related Work (Section~\ref{sec:related}) for this revision, leaving 14 Round~1 records in the primary corpus. Table~\ref{tab:prisma} summarizes the combined accounting.

\begin{table}[t]
\caption{Study selection process (PRISMA-style accounting).}
\label{tab:prisma}
\centering
\begin{tabular}{p{5.2cm}p{2.3cm}}
\hline
\textbf{Stage} & \textbf{Count} \\
\hline
Round 2 candidate records identified & 32 \\
Round 2 excluded (EX4/EX6: not a generated artifact / general GenAI-in-education, no specific system) & 11 \\
Round 2 excluded (EX3: general-purpose, non-education-specific) & 8 \\
Round 2 excluded (redundant with an included system) & 2 \\
Round 2 excluded (pre-dates LLM-based scope) & 1 \\
\textbf{Round 2 included} & \textbf{10} \\
Round 1 included (2026-08-06; candidate-record count not separately logged) & 15 \\
Round 1 reclassified to Related Work (EX5, meta-survey) & $-1$ \\
\hline
\textbf{Total included in primary corpus} & \textbf{24} \\
\hline
\end{tabular}
\end{table}

\subsection{Corpus Overview}

Of the 24 included primary studies, 21 are formatted as arXiv preprints (three of which are already accepted at a venue -- ICML 2026, EMNLP 2025, and CVPR 2025, respectively), 2 are CVPR 2024 proceedings papers (VBench, EvalCrafter), and 1 is a peer-reviewed journal article (Pellas~\cite{pellas2025aivideos}, \emph{Education Sciences}). This preprint-heavy composition is itself a property of a fast-moving field and is discussed as a threat to validity in Section~\ref{sec:threats}.

\section{Architectures: From Pixel Synthesis to Structured Intermediates}
\label{sec:architectures}

\subsection{The Decomposition Principle}

The architectural throughline across nearly every system in the corpus is decomposition. Rather than a single model mapping a prompt directly to pixels, each system inserts a checkable, editable, or executable intermediate representation between the language model and the rendered output, and evaluates or critiques that intermediate before rendering proceeds. This design choice recurs across otherwise very different systems and appears to be the field's dominant response to the correctness problem that end-to-end video diffusion cannot solve.

\subsection{Script and Event-Graph Intermediates}

Yan et al.\ propose LASEV, a hierarchical multi-agent system built around an Orchestrating Agent that supervises three specialized subordinates -- a Solution Agent, an Illustration Agent, and a Narration Agent \cite{yan2026lasev}. Rather than synthesizing pixels directly, LASEV constructs ``a structured executable video script'' that is deterministically compiled into the final video. The authors motivate this design explicitly by citing the failure of end-to-end video models on ``scenarios that require strict logical rigor and precise knowledge representation,'' and report production at over one million videos per day with more than a 95\% cost reduction relative to industry-standard approaches. These figures speak to scale and cost efficiency; as Section~\ref{sec:gaps} discusses, they say nothing about whether the resulting videos teach effectively.

\subsection{Code Intermediates}

Chen, Lin, and Shou take decomposition a step further with Code2Video, which uses a Planner--Coder--Critic agent pipeline to produce \emph{executable code} (Manim, a Python animation library) rather than a script or event graph \cite{chen2025code2video}. The Planner structures lecture content into a temporally coherent flow and prepares visual assets; the Coder converts that structure into executable Python with a scope-guided auto-fix loop; and the Critic, a vision-language model conditioned on ``visual anchor prompts,'' refines spatial layout before final rendering. Code2Video reports a 40\% improvement over direct code generation without the Planner--Critic loop, evaluated on a new benchmark (MMMC) discussed in Section~\ref{sec:benchmarks}.

\subsection{Multimodal Channel Decomposition}

A third decomposition strategy splits the output into independent, separately generated and separately alignable channels. Zhu, Lin, and Shou's Paper2Video introduces PaperTalker, a multi-agent framework that generates slides, subtitles, synthesized speech, and talking-head video as coordinated but distinct channels, with slide-wise generation parallelized for efficiency \cite{zhu2025paper2video}. Islam, Manik, and Wang's ALIVE follows a related pattern for live lecture delivery, combining avatar-delivered explanation \emph{as a pre-rendered lecture artifact} with a content-aware retrieval mechanism that fuses semantic similarity with timestamp alignment, and a real-time question-answering channel layered on top, all designed to run on local hardware for classroom deployment \cite{islam2025alive} -- the pre-rendered lecture component is what places it within IN1, distinct from the live-only agents excluded under EX4 (Table~\ref{tab:selection-criteria}). Holmberg's system for generating narrated lecture videos from existing slide decks similarly treats narration and visual highlighting as separate, alignable channels, using a highlight-alignment module that combines Levenshtein-distance matching with LLM-based semantic alignment at configurable granularity \cite{holmberg2025lectures}. Jo, Zhao, Liu, and Suzuki's Generative Lecture extends this family in a different direction -- rather than generating a lecture video from scratch, it converts an \emph{existing} lecture video into an interactive experience via an AI-cloned instructor (combining an avatar, synthesized voice, and an LLM) supporting on-demand clarification, adaptive quizzes, and personalized explanation \cite{jo2025generativelecture}.

\subsection{Slide-Centric Pipelines}

A large cluster of systems treats the slide deck itself as the primary generated artifact, with video (if produced at all) as a downstream rendering step. Zheng et al.'s PPTAgent uses a two-stage, edit-based methodology: it first analyzes reference presentations to extract slide-level functional types and content schemas, then drafts an outline and iteratively generates editing actions rather than pixels, evaluated with a companion framework, PPTEval, that scores Content, Design, and Coherence \cite{zheng2025pptagent}. Ge et al.'s AutoPresent, an eight-billion-parameter Llama-based model, generates slides as executable code and is reported to achieve results comparable to GPT-4o on this task, reinforcing the pattern -- also seen in Code2Video -- that programmatic generation outperforms direct image generation for structured visual content \cite{autopresent2025}. Aggarwal and Bhand's PASS extends this line of work beyond academic papers to arbitrary source documents, producing both slides and synthetic narration from an arbitrary Word document \cite{aggarwal2025pass}.

\subsection{The Animation-Specific Thread}
\label{sec:animation-thread}

A cluster of systems identified in the second search round applies the same decomposition principle specifically to \emph{animation} generation -- narrower in output modality than a full video pipeline, but architecturally the same wager: an intermediate representation that can be checked before rendering. Manim, an open-source Python animation library, has emerged as the dominant rendering substrate for this thread, extending Code2Video's use of it beyond a single system.

Samarth, Jain, Golugula, and Sathvik's Manimator was an early (mid-2025) instance of a two-stage pipeline converting a research paper or natural-language prompt into a structured scene description and then into Manim code \cite{samarth2025manimator}. Silva, Lotfi, Ihianle, Shahtahmassebi, and Bird address a question orthogonal to pipeline orchestration -- \emph{how the coder model itself should be trained} -- introducing ManimTrainer (supervised fine-tuning combined with reinforcement learning) and a renderer-in-the-loop inference strategy they also call ManimAgent, evaluated across 17 open-source sub-30B models \cite{silva2026manimtrainer}. Jiang et al.\ introduce a \emph{different}, unrelated system that also names itself ManimAgent (project name Paper2Manim), which instead of a fixed training procedure carries reflection experience across tasks via a dual-channel Episodic Memory Bank of success rationales and validated failure patterns, built entirely from its own task stream with no weight updates \cite{jiang2026manimagent} -- we flag this name collision explicitly, since both systems are legitimate, independently verified 2026 contributions.

Li et al.'s OmniManim targets a specific failure mode the training- and orchestration-focused systems above do not evaluate directly: visual defects (element overlap, misalignment) that only become apparent after rendering. Their Vision Agent predicts sparse keyframe layouts \emph{before} code generation, formalizing the task as ``render-feedback-aware constrained code generation'' \cite{li2026omnimanim}. Liao et al.'s ALGOGEN pushes the decomposition principle one level further for the algorithm-visualization sub-domain: rather than having an LLM generate animation code end to end, it decouples algorithm execution from rendering entirely, representing visual state as a ``Visualization Trace Algebra'' (VTA) that a deterministic renderer compiles via a ``Rendering Style Language'' (RSL) \cite{liao2026algogen}. Because the intermediate is a verifiable execution trace rather than free-form code, ALGOGEN's success rate is substantially higher than end-to-end baselines (99.8\% vs.\ 82.5\%) -- the clearest algorithm-visualization-specific demonstration of this survey's central architectural claim.

Two further systems in this thread take the animation-generation task in directions distinct from the training- and layout-focused systems above. Ku et al.'s TheoremExplainAgent is a second Manim-based system, but is notable for a second-order use of its own output: it generates long-form (5+ minute) videos explaining mathematical theorems, and uses the generated video to probe whether the \emph{generating} LLM's reasoning holds up under a more demanding, multimodal explanation task, not only to teach a human viewer \cite{ku2025theoremexplainagent}. Prasad and Mahapatra's whiteboard-generation system is the one member of this thread that departs from Manim entirely, replacing the code intermediate with a stroke-sequence representation: a vision-language model (Qwen2-VL-7B, fine-tuned via LoRA) predicts freehand drawing strokes synchronized to narrated speech, trained on the first available paired dataset of Excalidraw demonstrations with millisecond-precision timestamps \cite{prasad2026whiteboard}. This is architecturally the most distinct member of the animation thread: the checkable intermediate is neither code nor a script, but a structured drawing representation.

\subsection{The Diminishing Role of General-Purpose Text-to-Video Backbones}

General-purpose diffusion and transformer-based text-to-video models -- Sora, VideoPoet, CogVideoX, and open pipelines such as ModelScope and VideoCrafter2 -- remain active areas of research in their own right, and agentic general-purpose systems continue to appear (e.g.\ Huang et al.'s ViMax, which coordinates hierarchical narrative planning and visual-consistency tracking for long-form storytelling \cite{huang2026vimax}, screened but excluded under EX3 as non-education-specific). In some education-adjacent pipelines a general-purpose model still serves as the final rendering substrate. However, essentially none of the education-specific systems reviewed above treats such a model as the \emph{entire} pipeline. This is the paper's first architectural claim: agentic decomposition, not end-to-end pixel generation, is now the dominant architecture for LLM-based educational video and animation generation, precisely because decomposition is what makes intermediate correctness checking possible -- and the animation-specific thread in Section~\ref{sec:animation-thread} shows this principle generalizing to code, verifiable traces, and stroke sequences alike, not only to script- or slide-based intermediates.

\section{Benchmarks and Datasets}
\label{sec:benchmarks}

\subsection{General-Purpose Video-Generation Benchmarks}

Two benchmarks dominate general text-to-video evaluation. VBench, introduced by Huang et al.\ at CVPR 2024, decomposes ``video generation quality'' into sixteen hierarchical, disentangled dimensions -- among them subject consistency, motion smoothness, temporal flickering, and spatial relationships -- each with tailored prompts and automated evaluators validated against human perceptual judgments \cite{huang2024vbench}; an extended version, VBench++, broadens this evaluation suite further \cite{huang2024vbenchpp}. EvalCrafter, also published at CVPR 2024, takes a complementary approach, combining seventeen objective metrics spanning visual quality, content, and motion with subjective human ratings to produce a reliable cross-model ranking \cite{liu2024evalcrafter}. Both benchmarks are methodologically mature and widely adopted for their stated purpose. Neither, however, includes a dimension that measures the factual correctness of spoken or depicted content, or a viewer's comprehension of it -- an omission that is unsurprising given both were designed for general-purpose video generation, but one that matters a great deal once these benchmarks are used, as they routinely are, to justify claims about educational systems built on the same backbones.

\subsection{Education-Specific Benchmarks}

A second, much younger generation of benchmarks has emerged specifically to address this gap. SlidesBench, introduced alongside AutoPresent, comprises roughly seven thousand training and 585 testing examples drawn from 310 real slide decks across ten domains, and supports both reference-based evaluation (similarity to a target deck) and reference-free evaluation (design quality judged on its own terms) \cite{autopresent2025}. Code2Video introduces MMMC, a benchmark of professionally produced, discipline-specific educational videos, evaluated along VLM-as-judge aesthetic scores, code efficiency, and a novel metric called TeachQuiz \cite{chen2025code2video}. TeachQuiz works by having a vision-language model ``unlearn'' a concept and then measuring how much of that knowledge it recovers after watching the generated video -- the closest analogue in this literature to a genuine knowledge-transfer metric. Paper2Video introduces the largest and most metric-diverse benchmark found in this survey: 101 research papers paired with author-created presentation videos, slides, and speaker metadata, evaluated with four purpose-built metrics -- Meta Similarity, PresentArena, PresentQuiz, and IP Memory -- the last of which specifically measures whether a generated video preserves a paper's key claims \cite{zhu2025paper2video}. PresentQuiz follows the same architectural pattern as TeachQuiz: a quiz-based, model-judged proxy for whether the video conveys its intended content.

\subsection{Animation-Specific Benchmarks}

The animation-specific thread (Section~\ref{sec:animation-thread}) contributes its own generation of benchmarks, all introduced within the same search window as the systems they evaluate. Silva et al.\ introduce ManimBench alongside ManimTrainer/ManimAgent, evaluating 17 open-source models on Render Success Rate and Visual Similarity to a reference \cite{silva2026manimtrainer}. Li et al.\ construct two datasets for OmniManim -- ManimLayout-1K and EduRequire-500, the latter used as the evaluation benchmark -- targeting layout-quality metrics that neither ManimBench nor MMMC evaluates explicitly \cite{li2026omnimanim}. Liao et al.\ evaluate ALGOGEN on a 200-task LeetCode algorithm-visualization benchmark, reporting success rate as a correctness-style metric closer to code-execution verification than to visual-quality judgment \cite{liao2026algogen}. Ku et al.\ introduce TheoremExplainBench, 240 theorems across multiple STEM disciplines with five automated evaluation metrics, notable for evaluating multimodal video explanations against the reasoning of the generating model itself rather than only against a human audience \cite{ku2025theoremexplainagent}. Prasad and Mahapatra's whiteboard-generation dataset, while smaller (24 paired demonstrations across 8 STEM domains), is the first of its kind for this modality and evaluates temporal alignment between stroke generation and narration rather than any visual-quality or knowledge-transfer dimension \cite{prasad2026whiteboard}.

\subsection{The Meta-Survey Perspective}

Liu, Xiang, Li et al.'s AIGVE survey, discussed in Section~\ref{sec:related} as a prior secondary study rather than a primary study in this corpus, formally establishes ``AI-Generated Video Evaluation'' as a distinct research problem and explicitly argues that current frameworks remain insufficient for video's complex spatial and temporal dynamics \cite{liu2024aigve}. As noted in Section~\ref{sec:related}, this SLR's findings, detailed next, narrow that broader AIGVE critique specifically to pedagogical effectiveness.

\section{Evaluation Gaps}
\label{sec:gaps}

\subsection{What Is Measured, Ranked by Frequency}

Table~\ref{tab:eval-taxonomy} ranks the evaluation approaches found across the 24 primary studies in this corpus, from most to least common. Visual and production quality (VBench, EvalCrafter, ManimBench's Visual Similarity, the design dimension of PPTAgent's PPTEval) is measured almost universally. Structural and design coherence and textual similarity to a reference are also common. Knowledge-transfer proxies via an LLM or VLM judge -- TeachQuiz and PresentQuiz -- are present but rare; Section~\ref{sec:gaps} (Unvalidated Knowledge-Transfer Proxies) discusses the validation gap both share in detail. Human quality or usability judgment that stops short of a controlled learning-outcome comparison is a small but growing category, contributed by the animation-specific thread (ManimAgent) and the broader lecture-video literature (Generative Lecture). Dedicated, reported factual-accuracy metrics applied to a system's own generated educational content are effectively absent from every generation-system paper surveyed -- though ALGOGEN's verifiable-trace architecture is worth noting as a structural mitigation for its specific sub-domain (algorithm visualization), even though it is not reported as a factuality \emph{metric} per se. Controlled human learning-outcome studies, measured against a comparison condition, appear in three papers -- up from one in the narrower first-round search. Real classroom deployment with human quality judgment appears in one additional paper. Accessibility -- captioning quality, multilingual fidelity, or evaluation across learners with differing needs -- does not appear as a dedicated evaluation axis in any paper surveyed.

\begin{table}[t]
\caption{What current evaluation practice measures, ranked by prevalence across the 24 primary studies surveyed}
\label{tab:eval-taxonomy}
\centering
\begin{tabular}{p{4.6cm}p{2.6cm}}
\hline
\textbf{Evaluation target} & \textbf{Prevalence} \\
\hline
Visual / production quality & Near-universal \\
Structural / design coherence & Common \\
Textual similarity to reference & Common \\
Knowledge-transfer proxy (LLM/VLM judge) & Rare (2 papers) \\
Human quality/usability judgment (non-learning-outcome) & Rare (2 papers) \\
Factual accuracy of generated content & Effectively absent \\
Controlled human learning-outcome study & 3 papers \\
Real classroom deployment + student judgment & 1 paper \\
Accessibility as a dedicated axis & Absent \\
\hline
\end{tabular}
\end{table}

\subsection{The Learning-Outcome Gap}

Three papers in this corpus report a genuine controlled comparison against a human learning outcome, and all three are worth examining closely precisely because they remain exceptions rather than the norm. Stavrinou et al.'s ReelsEd system converts long-form lecture video into short-form ``reels'' while preserving instructor-authored material, and the authors evaluate it with a controlled user study of 62 university students, comparing reels against traditional long-form video on engagement, quiz performance, task efficiency, and cognitive load \cite{stavrinou2025reeldeal}. The reels condition outperformed long-form video on engagement, quiz performance, and task efficiency without an increase in cognitive load. This remains, to our knowledge, the strongest learning-outcome evaluation in the current literature on LLM-based educational video generation.

Joshi et al.'s LLM2Manim is the first learning-outcome study specific to \emph{generated animation} rather than video or slides: a within-subject A-B study with 100 undergraduate students found modestly higher post-test scores for animation-based instruction (83\% vs.\ 78\%, $p<.001$ -- a gap the authors themselves characterize as ``slight''), alongside a medium-to-large learning-gains effect ($d=0.67$), a large engagement effect ($d=0.94$), and a small-to-medium reduction in cognitive load ($d=0.41$) \cite{joshi2026llm2manim}. Pellas's study of AI-generated instructional video in science teacher education is the only journal-published (rather than preprint) human-subject study in this corpus, and the only one conducted with a real classroom-adjacent population (55 Greek pre-service science teachers): both video formats tested (with and without an embedded preview feature) effectively supported self-efficacy, task performance, and knowledge retention, with no significant difference between the two conditions \cite{pellas2025aivideos}.

Leinonen, Zhang, and Hellas take a complementary approach, deploying AI-generated slides -- built from instructor notes using five different tools and pre-screened for quality by instructors -- in real university courses, then asking students to blind-rate the deployed slides against instructor-created ones \cite{leinonen2026aislides}. Students could not reliably distinguish AI-generated slides from instructor-made ones and rated them similarly in overall quality, though slides students merely \emph{suspected} of being AI-generated were rated lower regardless of their actual origin -- a perception bias that no artifact-only benchmark could ever detect, since it depends entirely on a human viewer's beliefs about provenance rather than any property of the artifact itself.

A further two papers report human evaluation that stops short of a controlled learning-outcome comparison. Jiang et al.'s ManimAgent reports blind human Pass@1 judgments of rendered animation quality as memory size grows \cite{jiang2026manimagent}, and Jo et al.'s Generative Lecture reports a small design-elicitation study (N=8) and a separate usability study (N=12) of an interactive lecture-augmentation feature set \cite{jo2025generativelecture}. Neither compares a generated artifact against a control condition on a learning outcome; both are nonetheless evidence that human evaluation, in some form, is becoming somewhat more common in the newer wave of systems than in the video-pipeline-centric literature that Round 1 of this review originally covered.

Taken together, six of the 24 primary studies (25\%) report some form of human evaluation, but only three (12.5\%) report a controlled learning-outcome comparison, and one further study (4\%) reports real classroom deployment with human quality judgment. This is a meaningful increase over the two-of-nineteen figure reported in the narrower first-round search, and is driven specifically by the animation-specific thread: LLM2Manim, the animation-specific controlled learning-outcome study, was found only once the search was widened beyond video pipelines in Round 2. Pellas's study, while about instructional video rather than animation, was likewise a Round 2 addition, since Round 1's search did not surface it. Even so, at roughly one in eight studies, the field's evaluation practice still overwhelmingly measures the generated artifact rather than the learner -- the central finding of this review is updated, not overturned, by the wider search.

\subsection{Unvalidated Knowledge-Transfer Proxies}

TeachQuiz and PresentQuiz remain the most promising evaluation innovations to emerge from this literature, precisely because they attempt to measure something closer to comprehension than to visual fidelity. Both work by having a language or vision-language model answer quiz questions after processing the generated video, and both are reasonable, scalable engineering solutions to the problem that human-subject studies are slow and expensive to run at the scale needed to iterate on a generation system. Neither paper that introduces one of these metrics, however, reports a study correlating the model-judge score with actual human quiz performance on the same video set. This is a specific, tractable, and currently open validation gap: without it, a system could in principle improve its TeachQuiz or PresentQuiz score by becoming better at satisfying a language model's quiz-answering behavior specifically, rather than by becoming more pedagogically effective for a human audience -- and the current literature has no evidence to rule that possibility out.

\subsection{Factual Accuracy: A Known Problem, Nowhere Directly Measured}

A substantial and separate literature documents that LLMs cannot reliably self-evaluate the factual accuracy of their own output without external grounding such as retrieval-augmented generation, knowledge graphs, or human oversight \cite{hallucination2026review}. This is directly relevant to educational content, where a confidently delivered factual error is pedagogically worse than a stylistic or design flaw, since a learner has no independent way to catch it. Despite this, none of the generation systems surveyed in Section~\ref{sec:architectures} reports a dedicated fact-checking pass, applied to its own generated narration or slide text, as part of its evaluation. The closest structural analogue is ALGOGEN's verifiable-trace architecture (Section~\ref{sec:animation-thread}), which sidesteps rather than measures the problem: by construction, its rendered output cannot diverge from the underlying algorithm trace, for the narrow sub-domain of algorithm visualization. This is a promising direction but does not generalize to open-ended narrated content, where applying an existing fact-checking or hallucination-detection method and reporting the result remains a direct, comparatively low-effort extension available to any of the systems reviewed here.

\subsection{Accessibility: An Absent Axis}

Several of the systems surveyed -- ALIVE, ReelsEd, and Generative Lecture in particular -- are explicitly designed for real classroom or self-study use, yet no paper in this survey reports caption quality, multilingual generation fidelity, or evaluation across learners with differing accessibility needs (for example, cognitive load for neurodivergent learners, or comprehension for non-native speakers of the narration language) as a dedicated evaluation axis. Given the stated deployment goals of these systems, this is a significant and directly addressable gap.

\section{Threats to Validity}
\label{sec:threats}

\textbf{Single-reviewer selection.} This SLR was conducted and screened by a single reviewer rather than two or more independent screeners reconciling disagreements, unlike the reference SLR methodology this review otherwise follows. Individual title/abstract screening judgments -- particularly borderline EX3/EX6 calls on systems with partial educational framing -- may not replicate under a second, independent screener.

\textbf{Search method, not database querying.} Both rounds used web-based search engines against arXiv, ACL Anthology, CVF Open Access, and general indexed sources, rather than direct Boolean queries against subscription bibliographic databases (IEEE Xplore, ACM Digital Library, Scopus, Web of Science). This is a materially different search method from the reference SLR this paper's Methodology structure is modeled on, and likely under-covers venue-gated or non-preprint work relative to a multi-database search. It is a defensible choice given this field's arXiv-first publication norms (Section~\ref{sec:methodology} notes 20 of 24 included studies are preprints), but it is not equivalent to a database-indexed search, and should be weighed accordingly by any reader using this review's corpus as a completeness claim rather than a representative sample.

\textbf{Round 1 record-level accounting was not preserved.} Table~\ref{tab:prisma}'s Round 1 row reports only the final included count (15), not the number of candidate records screened to reach it, because that count was not logged at the time. This is an honest gap in this review's own auditability and is disclosed rather than backfilled with an estimate.

\textbf{Publication and availability bias.} The corpus is heavily arXiv-preprint-weighted (Section~\ref{sec:methodology}), which means findings, benchmark results, and reported effect sizes have not, in most cases, passed peer review. Preprint results in a fast-moving field are also subject to revision between versions; this review captured a snapshot as of 2026-08-13 and does not track subsequent revisions.

\textbf{Recency and velocity risk.} New architectures and benchmarks in this space have appeared on a roughly monthly cadence over the period this review covers (Section~\ref{sec:discussion}). A review of a fast-moving field risks being stale within months of completion; this is disclosed rather than mitigated, since no search methodology available to a single reviewer can outpace the publication rate itself.

\textbf{Terminology inconsistency.} The corpus includes two unrelated systems that both name themselves ``ManimAgent'' (Section~\ref{sec:animation-thread}). This is disclosed explicitly at the point of citation to avoid reader confusion, but it is representative of a broader pattern in a fast-moving, preprint-heavy field: system names are not yet standardized, which complicates both literature search (a name-based search for one system will also surface its unrelated namesake) and citation tracking.

\section{Discussion and Future Directions}
\label{sec:discussion}

The architecture and benchmark literatures surveyed in Sections~\ref{sec:architectures} and~\ref{sec:benchmarks} are converging rapidly: new agentic architectures and new education- or animation-specific benchmarks have appeared on a roughly monthly cadence over the period covered by this review, including an entire animation-specific thread that this review's own first round missed. The evaluation-methodology literature has not kept pace. The AIGVE meta-survey's call for more sophisticated evaluation frameworks \cite{liu2024aigve} has not yet produced anything for pedagogical effectiveness comparable to what VBench achieved for general video quality \cite{huang2024vbench} -- a single, widely adopted, multi-dimensional standard that most new systems are evaluated against by default. We see this as the field's most obviously missing artifact: a benchmark for pedagogical effectiveness that plays the same role for education- and animation-specific generation that VBench plays for general video quality.

Five concrete next steps follow directly from the gaps identified in Section~\ref{sec:gaps}, updated from the four identified in the narrower first-round search. First, and most urgently: a cross-system, controlled human-subject study. Every learning-outcome result currently in the literature comes from a single research group evaluating its own single system in isolation; running the same protocol used by Stavrinou et al.\ \cite{stavrinou2025reeldeal} across two or three systems from different research groups and modalities -- for instance Code2Video, LLM2Manim, and a general text-to-video baseline -- would be the field's first genuinely comparative, cross-modality learning-outcome result. Second, TeachQuiz and PresentQuiz should be validated against human quiz performance on a shared set of generated videos; this is a focused, tractable study that would substantially increase confidence in the field's most promising existing proxy metrics, or reveal that they need revision. Third, applying an existing factuality or hallucination-detection method to the generated narration or slide text of one of the systems reviewed here, and reporting the result, would close a gap that is currently open purely for want of anyone doing it, not for lack of available technique -- ALGOGEN's verifiable-trace approach (Section~\ref{sec:animation-thread}) is a template worth generalizing beyond algorithm visualization specifically. Fourth, accessibility evaluation -- caption accuracy, multilingual fidelity, and comprehension across learners with differing needs -- deserves the same systematic treatment currently given to visual quality, particularly for the systems in this survey that are explicitly aimed at real classroom deployment. Fifth, a multi-reviewer, multi-database replication of this SLR -- addressing the single-reviewer and search-method limitations disclosed in Section~\ref{sec:threats} -- would materially strengthen confidence in both the architecture taxonomy and the evaluation-gap quantification that this review reports.

\section{Conclusion}
\label{sec:conclusion}

LLM-based educational video and animation generation has, in a short period, converged on a clear architectural answer: decompose the generation task into a structured, checkable intermediate -- a script, an event graph, executable code, a verifiable trace, a stroke sequence, or an editable slide deck -- rather than generating pixels end to end. This decomposition is a direct and, we argue, correct response to the specific failure mode that pixel-space video diffusion cannot address: guaranteeing that educational content is factually and logically sound. This systematic review's animation-specific thread shows the same principle generalizing cleanly to a modality -- narrated STEM animation, algorithm visualization, whiteboard drawing -- that a narrower, video-only search would miss entirely. Benchmarking has responded in kind, producing a new generation of education- and animation-specific evaluation suites that represent genuine progress over general-purpose video-quality metrics. What the field has not yet produced, with only three controlled learning-outcome studies and one classroom-deployment study among the 24 primary studies reviewed here, is broad evidence that these systems actually improve human learning. Closing that gap -- through cross-system, cross-modality human-subject studies, validated knowledge-transfer proxies, dedicated factuality checks, accessibility evaluation, and a multi-reviewer replication of this SLR itself -- is, in our view, the highest-leverage open problem in this area, and the natural next chapter for a field that has so far mostly asked whether its output looks right rather than whether it teaches.

\bibliographystyle{IEEEtran}
\bibliography{../Bibliography_base}

\end{document}
